\documentclass[12pt]{article}
\usepackage[utf8]{inputenc}
\usepackage[T1]{fontenc}
\usepackage{hyperref}
\title{Wagawaga Language Resources for K-12 Curriculum Development}
\author{Compiled by Claude Sonnet 5.5 (claude-sonnet-5-5) for Andr\'as Kornai}
\date{}
\begin{document}
\maketitle

\section{Introduction}
The code wgw is a retired ISO 639-3 code for Wagawaga (Glottocode waga1262). The Glottolog family label ``Bookkeeping'' marks it as an administrative entry, not a genealogical family. The speech variety is Wagawaga, an Oceanic (Austronesian) language of Milne Bay Province, Papua New Guinea. Sources attach it to the active codes wgb (Wagawaga, also Baeaula) and ylb (Yaleba, also Gamadoudou). Wikipedia says the Gamadoudou, Soma'a and Sileba dialects may form a separate language, Yaleba. The same Wikipedia article gives about 1,900 native speakers (2000 census) and is flagged as relying largely on a single source. Very little educational material was found.

\section{Primary Educational Resources}
No primer, textbook or school programme for Wagawaga turned up in the searches. Scripture Earth carries a listing under code wgb, but Joshua Project reports Bible translation as only started, with no portions or New Testament reported.

\section{Linguistic and Reference Materials}
\begin{itemize}
\item Wikipedia, Wagawaga language (New Guinea): \url{https://en.wikipedia.org/wiki/Wagawaga_language_(New_Guinea)}
\item Rosetta Project entry for wgw, with alternate names Gamadoudou, Kilakilana, Baeaula and Gibara: \url{https://rosettapanglossia.longnow.org/index.php?title=Wgw}
\item A search result mentions Arthur Capell's Yaleba materials held at Paradisec. I did not see a direct URL for them and did not check their contents, so they are not listed with a link.
\end{itemize}
No dictionary or wordlist was found.

\section{Religious and Community Resources}
\begin{itemize}
\item Joshua Project, Wagawaga/Deamuni people profile: \url{https://joshuaproject.net/people_groups/15816}
\item Joshua Project, language page (wgb): \url{https://joshuaproject.net/index.php/languages/wgb}
\item Scripture Earth listing for wgb: \url{https://www.scriptureearth.org/00eng.php?iso=wgb}
\end{itemize}

\section{Digital and Audio Resources}
\begin{itemize}
\item Global Recordings Network, Wagawaga audio Bible stories and evangelistic messages: \url{https://globalrecordings.net/en/language/wgb}
\end{itemize}

\section{Curriculum Development Context}
Wagawaga is not known to be taught in schools, and there are no teaching materials. A curriculum would have to start from new documentation with the community. Work on the dialect question (Wagawaga proper versus Yaleba) would need to come first. The English-medium material above gives only background.
\end{document}